% Standard pen-and-paper worksheet preamble, in the form m01 and m02 use.
% Compile with: xelatex -interaction=nonstopmode exercise.tex   (run twice)
%
% 14pt, not 12pt: the sheet is read with a pencil in hand. There are no blanks
% and no ruled answer gaps anywhere on it -- questions are asked in words and
% answered on the back -- so the paper those would have cost goes to the text
% and to the drawings.
\documentclass[a4paper, 14pt]{extarticle}
\usepackage[top=0.75in, bottom=0.7in, left=0.85in, right=0.85in]{geometry}
\usepackage{amsmath}
\usepackage{amssymb}
\usepackage{array}
\usepackage{graphicx}
\usepackage[hidelinks]{hyperref}   % hidelinks, or the print carries colour boxes
\usepackage{tcolorbox}
\usepackage{enumitem}
\usepackage{etoolbox}
\usepackage{tikz}
\usepackage{fontspec}
\usetikzlibrary{calc,decorations,patterns,arrows,decorations.pathmorphing,positioning}
\definecolor{pltblue}{HTML}{1F77B4}

% Body face: Charter, and the four ways to get it. XCharter *is* Charter and
% ships with a full TeX Live, so it is asked for first and by file -- a CI
% runner answers yes to \IfFontExistsTF{Charter}, then cannot typeset it, which
% is how these sheets stopped building there.
\IfFontExistsTF{XCharter-Roman.otf}{
  \setmainfont{XCharter}[
    Extension      = .otf,
    UprightFont    = *-Roman,
    BoldFont       = *-Bold,
    ItalicFont     = *-Italic,
    BoldItalicFont = *-BoldItalic]
}{
  \IfFontExistsTF{Charter}{\setmainfont{Charter}}{
    \IfFontExistsTF{Iowan Old Style}{\setmainfont{Iowan Old Style}}{
      \IfFontExistsTF{Georgia}{\setmainfont{Georgia}}{}}}}

% Handwriting face, resolved once into \ppfont. Used for a heading or two (the
% answer key's "Answers"), not for the body and not for the drawings: the
% figures are read closely and want the text face.
%
% Naming a font that is not installed is a hard error, not a fallback, which is
% why this is a chain and not a line: seven sheets used to name one outright and
% stopped building the day it was uninstalled. Write \ppfont, never
% \fontspec{<a name>}. Excalifont is vendored in tools/fonts/ and installed by
% CI, so it is the one face that is always present.
\newcommand{\ppfont}{}
\IfFontExistsTF{Pretty Neat}{\renewcommand{\ppfont}{\fontspec{Pretty Neat}}}{
  \IfFontExistsTF{Humor Sans}{\renewcommand{\ppfont}{\fontspec{Humor Sans}}}{
    \IfFontExistsTF{Excalifont}{\renewcommand{\ppfont}{\fontspec{Excalifont}}}{
      \renewcommand{\ppfont}{}}}}

\setlength{\parindent}{0pt}
\setlength{\parskip}{0.38em}
\setlength{\emergencystretch}{1.2em}
\raggedbottom

% Break the page unless #1 of it is left.
%
% This is needspace's glue trick, written out because needspace.sty lives in
% texlive-latex-extra and not in a BasicTeX install. Do NOT rewrite it as a test
% on \pagegoal minus \pagetotal: the page builder has not necessarily run when
% the macro expands, so those numbers can still be the previous page's, the test
% says "no room" on a page that is empty, and the \newpage throws that page
% away. m02 shipped a blank page 3 that way, and it is invisible in the source.
\newcommand{\needroom}[1]{%
  \par
  \begingroup
    \setlength{\dimen0}{#1}%
    \vskip 0pt plus \dimen0
    \penalty -100%
    \vskip 0pt plus -\dimen0
  \endgroup}

% A part heading. The sheet has no title: it opens on Part 1. The heading
% carries the break test, because a heading printed at the foot of a page with
% nothing under it is worse than a short page.
\newcommand{\parthead}[2]{%
  \needroom{6\baselineskip}%
  \partheadhere{#1}{#2}}

% The same heading with no break test, for one already inside a minipage --
% \newpage means nothing in there.
\newcommand{\partheadhere}[2]{%
  \vspace{0.5em}
  {\large\bfseries Part #1 --- #2}\par
  \vspace{-0.4em}\rule{\textwidth}{0.8pt}\par
  \vspace{0.1em}}

% Centred fixed-width cell, for tables the student writes into.
\newcolumntype{C}[1]{>{\centering\arraybackslash}p{#1}}

% xkcd-style hand-drawn line decoration.
% WARNING: applying [xkcd] to many nodes can throw "Dimension too large".
% Use it on a handful of \draw lines only, never on every node.
\makeatletter
\pgfset{
  /pgf/decoration/randomness/.initial=2,
  /pgf/decoration/wavelength/.initial=100
}
\pgfdeclaredecoration{sketch}{init}{
  \state{init}[width=0pt,next state=draw,persistent precomputation={
    \pgfmathsetmacro\pgf@lib@dec@sketch@t0
  }]{}
  \state{draw}[width=\pgfdecorationsegmentlength,
  auto corner on length=\pgfdecorationsegmentlength,
  persistent precomputation={
    \pgfmathsetmacro\pgf@lib@dec@sketch@t{mod(\pgf@lib@dec@sketch@t+pow(\pgfkeysvalueof{/pgf/decoration/randomness},rand),\pgfkeysvalueof{/pgf/decoration/wavelength})}
  }]{
    \pgfmathparse{sin(2*\pgf@lib@dec@sketch@t*pi/\pgfkeysvalueof{/pgf/decoration/wavelength} r)}
    \pgfpathlineto{\pgfqpoint{\pgfdecorationsegmentlength}{\pgfmathresult\pgfdecorationsegmentamplitude}}
  }
  \state{final}{}
}
\tikzset{xkcd/.style={decorate,decoration={sketch,segment length=0.5pt,amplitude=0.5pt}}}
\makeatother

% Figures shared by exercise.tex and solutions.tex. Edit here, never in a copy.
%
% The nine stations, in fixed positions. #1 spreads them: it is the width in cm
% of one layout unit, so the circles keep the same size on every copy and only
% the room to draw lines between them changes.
% Nodes are named A..I so an edge can be drawn on top of the layout.
\newcommand{\gridnodes}[1]{%
  \begin{tikzpicture}[x=#1 cm, y=#1 cm]
    \node[draw, circle, inner sep=0.12cm] (I) at (0.5, 1)    {I};
    \node[draw, circle, inner sep=0.12cm] (D) at (3, 1)       {D};
    \node[draw, circle, inner sep=0.12cm] (C) at (4.2, -0.2)  {C};
    \node[draw, circle, inner sep=0.12cm] (B) at (2, -0.5)    {B};
    \node[draw, circle, inner sep=0.12cm] (E) at (0.5, -1.6)  {E};
    \node[draw, circle, inner sep=0.12cm] (A) at (3, -2.5)    {A};
    \node[draw, circle, inner sep=0.12cm] (F) at (-0.5, -0.5) {F};
    \node[draw, circle, inner sep=0.12cm] (H) at (1, -3)      {H};
    \node[draw, circle, inner sep=0.12cm] (G) at (4.5, -2)    {G};
  \end{tikzpicture}}

% The cost table.
\newcommand{\costtable}[1]{%
  \scalebox{#1}{%
    \begin{tabular}{|c|c|c|c|c|c|c|c|c|c|}
      \hline
         & A  & B  & C  & D  & E  & F  & G  & H  & I  \\ \hline
      A  & 0  & 2  & 8  & 12 & 4  & 10 & 14 & 6  & 16 \\ \hline
      B  & -  & 0  & 3  & 2  & 1  & 5  & 11 & 15 & 7  \\ \hline
      C  & -  & -  & 0  & 4  & 10 & 14 & 6  & 12 & 16 \\ \hline
      D  & -  & -  & -  & 0  & 5  & 11 & 15 & 7  & 13 \\ \hline
      E  & -  & -  & -  & -  & 0  & 6  & 12 & 16 & 8  \\ \hline
      F  & -  & -  & -  & -  & -  & 0  & 7  & 13 & 9  \\ \hline
      G  & -  & -  & -  & -  & -  & -  & 0  & 8  & 10 \\ \hline
      H  & -  & -  & -  & -  & -  & -  & -  & 0  & 11 \\ \hline
      I  & -  & -  & -  & -  & -  & -  & -  & -  & 0  \\ \hline
    \end{tabular}}}

% The empty plotting grid for the two curves.
% x: number of stations removed, 0..9.  y: connectivity, 0..1.
\newcommand{\profilegrid}{%
  \begin{tikzpicture}[x=1.7cm, y=9.6cm]
    \draw[step=1, gray!30, very thin] (0,0) grid[xstep=1, ystep=0.1] (9,1);
    \draw[-latex, thick] (0,0) -- (9.4,0);
    \draw[-latex, thick] (0,0) -- (0,1.08);
    \foreach \x in {0,1,2,3,4,5,6,7,8,9}
      \draw (\x,0) -- (\x,-0.015) node[below, font=\small] {\x};
    \foreach \y/\l in {0/0, 0.2/0.2, 0.4/0.4, 0.6/0.6, 0.8/0.8, 1/1}
      \draw (0,\y) -- (-0.08,\y) node[left, font=\small] {\l};
    \node[below, font=\small] at (4.5,-0.085) {stations removed};
    \node[rotate=90, font=\small, anchor=south] at (-0.62,0.5) {connectivity};
  \end{tikzpicture}}


\begin{document}

\parthead{1}{Nine stations and no money}

\begin{minipage}{0.52\textwidth}
    {\bf Question 1}:
    Let us design a small power grid network with nine stations A--I (nodes).
    Initially, these stations are isolated.
    The cost associated with each connection is provided in the table.
    Minimize the total cost of the lines to build a single connected network.

    Draw your lines on the nine stations below.
\end{minipage}
\hfill
\begin{minipage}{0.46\textwidth}
\begin{center}
    \costtable{0.82}
\end{center}
\end{minipage}

\begin{center}
  \gridnodes{2.55}
\end{center}

{\bf Question 2}:
Compare your grid with your neighbour's.
In what order did each of you add the lines?
Did the two of you end up paying the same total cost?

{\bf Question 3}:
Pick any two stations in your grid.
How many different routes run between them?

\parthead{2}{Break it}

\begin{minipage}{\textwidth}
{\bf Question 4}:
Consider the scenario where a single station fails and disconnects from the network.
Enumerate all possible cases, and compute the size of the largest connected component in each case.
Write the sizes in the middle column.

\begin{center}
\begin{tabular}{|C{3.6cm}|C{4.6cm}|C{4.2cm}|}
\hline
Station that fails & Size of the largest connected component & That size divided by 9 \\ \hline
A & & \\[0.76cm] \hline
B & & \\[0.76cm] \hline
C & & \\[0.76cm] \hline
D & & \\[0.76cm] \hline
E & & \\[0.76cm] \hline
F & & \\[0.76cm] \hline
G & & \\[0.76cm] \hline
H & & \\[0.76cm] \hline
I & & \\[0.76cm] \hline
\end{tabular}
\end{center}
\end{minipage}

{\bf Question 5}:
Divide each of those sizes by 9 and write the result in the last column.
That fraction is called the \textbf{connectivity} of the damaged grid.
On average, how much connectivity remains?

\needroom{5\baselineskip}

{\bf Question 6}:
A power-grid network is subjected to a \textit{targeted attack} that breaks the network by removing one node.
The attacker targets a single station whose removal minimizes the size of the largest connected component in the resulting network.
Which station is most susceptible to this attack?

\parthead{3}{One at a time}

\begin{minipage}{\textwidth}
{\bf Question 7}:
Now the stations go one after another, and two orders are printed below.
Take your grid from Question 1, remove the stations in the given order, and after every removal write the connectivity into the table.

\vspace{0.2em}
\begin{tabular}{@{}ll@{}}
\textbf{Order R} (an accident): & I, C, H, E, A, G, D, F, B \\[0.2em]
\textbf{Order T} (the attacker): & B, A, C, D, E, F, G, H, I \\
\end{tabular}
\vspace{0.3em}

\begin{center}
\small
\begin{tabular}{|l|C{1.0cm}|C{1.0cm}|C{1.0cm}|C{1.0cm}|C{1.0cm}|C{1.0cm}|C{1.0cm}|C{1.0cm}|C{1.0cm}|}
\hline
stations removed & 0 & 1 & 2 & 3 & 4 & 5 & 6 & 7 & 8 \\ \hline
Order R & 1 & & & & & & & & \\[0.3cm] \hline
Order T & 1 & & & & & & & & \\[0.3cm] \hline
\end{tabular}
\end{center}

Then plot both rows on the grid, and join each one into a curve.

\begin{center}
  \profilegrid
\end{center}
\end{minipage}

{\bf Question 8}:
The area under a curve is a single number for how much of the grid survived the whole sequence.
Estimate that area for each of your two curves.

{\bf Question 9}:
The same grid holds up under one of those two orders and falls apart under the other.
Say, in your own words, what makes it strong against one and weak against the other.

\parthead{4}{Build it back}

\begin{minipage}{\textwidth}
{\bf Question 10}:
Redesign the network so that it is as strong as possible against targeted attacks.
To simplify, ignore the cost of the lines but keep the same number of lines as your grid in Question 1.
You can change the pre-existing lines completely.
The attacker can remove only one station.

\begin{center}
  \gridnodes{1.8}
\end{center}
\end{minipage}

\begin{minipage}{\textwidth}
{\bf Question 11}:
Now the attacker can remove two stations at once, and you have additional budget to add four extra lines to protect the network.
Design a network of \underline{\textbf{12}} edges so that it is as strong as possible against the targeted attacks.

\begin{center}
  \gridnodes{1.8}
\end{center}
\end{minipage}

\end{document}
